\section{KITTI dataset}
\subsection{Dataset Information}
KITTI dataset~\cite{KITTI} is one of the most common datasets used in research varying from computer vision to robotics. It includes camera images, laser scans, high-precision GPS measure-
ments, and IMU accelerations from a combined GPS/IMU system, capturing real-world traffic situations and range from freeways over rural areas to inner-city scenes with many static and dynamic objects. Many SOTA research of Object Detection, Depth Estimation, Semantic Seg-mentation used KITTI as a base metric for comparison. \Cref{fig:kitti_example} illustrates a sample of images from KITTI dataset.

KITTI consists of 7481 training images and 7518 test images for each task. Despite the fact that there are eight different classes, only the classes \textit{Car, Cyclist,} and \text{Pedestrian} are evaluated in their benchmark, as only for those classes have enough instances for a comprehensive evaluation have been labeled. The dataset’s sub-folders are structured as follows:
\begin{itemize}
    \item \texttt{image\_02}: contains the left color camera images (png)
    \item \texttt{label\_02}: contains the left color camera label files (plain text files)
    \item  \texttt{calib}: contains the calibration for all four cameras (plain text file)
\end{itemize}

\begin{figure}
    \centering
    \includegraphics[width=1\linewidth]{images//dataset//kitti/kitti_example.png}
    \caption{Examples from the KITTI dataset (from left color camera).}
    \label{fig:kitti_example}
\end{figure}

The label files contain 15 columns for each object:
\begin{itemize}
    \item \texttt{type:} Describes the type of object: ‘Car’, ‘Van’, ‘Truck’, ‘Pedestrian’, ‘Person sitting’, ‘Cyclist’, ‘Tram’, ‘Misc’ or ‘DontCare’.
    \item \texttt{truncated:} Float from 0 (non-truncated) to 1 (truncated), where truncated refers to the object leaving image boundaries.
    \item \texttt{occluded:} Integer (0,1,2,3) indicating occlusion state: 0 means fully visible, 1 means partly occluded 2 means largely occluded and 3 means unknown.
    \item $\alpha:$ Observation angle of object, ranging in $\left[-\pi,\pi\right]$.
    \item \texttt{bounding box:} 2D bounding box of object in the image (0-based index): contains left, top, right, bottom pixel coordinates.
    \item \texttt{dimensions:} 3D object dimensions: height, width, length (in meters).
    \item \texttt{location:} 3D object location x,y,z in camera coordinates (in meters).
    \item $r_y:$ Rotation $r_y$ around Y-axis in camera coordinates, ranging in $\left[-\pi,\pi\right]$.
\end{itemize}
The coordinates in the camera coordinate system can be projected in the image by using the $3\times 4$ projection matrix in the calibration folder, where for the left color camera for which the
images are provided, P2 must be used~\cite{KITTI}. The difference between $r_y$ and $\alpha$ is that $r_y$ is directly
given in camera coordinates, while $\alpha$ also considers the vector from the camera center to the
object center, to compute the relative orientation of the object with respect to the camera.
\subsection{Evaluation Metrics}

The KITTI object detection benchmark evaluates performance using the Average Precision (AP) metric, following the PASCAL VOC standard. To determine if a prediction is a True Positive (TP), the Intersection over Union (IoU) between the predicted bounding box and the ground truth bounding box must exceed a specific threshold. 

For the \textit{Car} class, the required IoU overlap is $70\%$, while for the \textit{Pedestrian} and \textit{Cyclist} classes, the required IoU is $50\%$. Detections that do not meet these overlap criteria or are assigned to a ground truth that has already been detected are counted as False Positives (FP).

To allow for a nuanced analysis of detector performance across different scenarios, the evaluation is performed separately for three difficulty levels: \textit{Easy}, \textit{Moderate}, and \textit{Hard}. These levels are defined based on the minimum height of the bounding box, the occlusion level, and the truncation ratio, as detailed in \Cref{tab:kitti_difficulty}.

\begin{table}[h]
    \centering
    \caption{Definition of KITTI difficulty levels for object detection evaluation.}
    \label{tab:kitti_difficulty}
    \begin{tabular}{l|ccc}
        \hline
        \textbf{Difficulty} & \textbf{Min. Height} & \textbf{Max. Occlusion} & \textbf{Max. Truncation} \\
        \hline
        \textit{Easy} & 40 Px & Fully Visible (0) & 15\% \\
        \textit{Moderate} & 25 Px & Partly Occluded (1) & 30\% \\
        \textit{Hard} & 25 Px & Difficult to See (2) & 50\% \\
        \hline
    \end{tabular}
\end{table}

Historically, the KITTI benchmark utilized an 11-point interpolation for calculating $AP|_{R_{11}}$. However, current state-of-the-art research typically utilizes the newer 40-point interpolation metric $AP|_{R_{40}}$, introduced in~\cite{MonoDis} to resolve interpolation bias and provide a more accurate comparison. The final ranking of methods on the leaderboard is determined by the performance on the \textbf{Moderate} difficulty setting.